% ==============================================================
% Kapitel 5: Theoretische Grundlagen
% ==============================================================

\section{Theoretische Grundlagen}

\subsection{Containerisierung (Docker)}
Containerisierung ist eine Methode der Softwarebereitstellung, bei der eine 
Anwendung zusammen mit allen ihren Abhängigkeiten -- wie Laufzeitumgebungen, 
Bibliotheken und Konfigurationsdateien -- in einer isolierten Einheit, dem 
sogenannten Container, verpackt wird. Im Gegensatz zu herkömmlichen virtuellen 
Maschinen teilen sich Container den Kernel des Hostsystems und benötigen daher 
keine vollständige Betriebssysteminstallation. Dies macht sie deutlich 
ressourcenschonender und schneller startbar.

Docker ist die am weitesten verbreitete Plattform zur Erstellung und Verwaltung 
von Containern. Die Grundlage eines jeden Docker-Containers bildet ein 
\textit{Image} -- eine unveränderliche Vorlage, die mithilfe eines 
\textit{Dockerfiles} definiert wird. In diesem Dockerfile werden alle 
Schritte festgelegt, die notwendig sind, um die Anwendung in einen 
lauffähigen Zustand zu versetzen. Aus einem Image können beliebig viele 
Container instanziiert werden, die jeweils isoliert und reproduzierbar 
laufen.

Ein wesentlicher Vorteil von Docker ist die Plattformunabhängigkeit: Ein 
einmal erstelltes Image verhält sich auf jedem System identisch, unabhängig 
davon, ob es auf einem lokalen Entwicklungsrechner, einem Server oder einem 
Einplatinencomputer wie dem Raspberry Pi betrieben wird. Dies löst das 
klassische Problem \textit{„works on my machine"} und erleichtert sowohl 
die Entwicklung als auch den Betrieb erheblich.

Für die Verwaltung mehrerer Container, die gemeinsam eine Anwendung bilden, 
bietet Docker mit \textit{Docker Compose} ein Werkzeug zur Definition und 
Orchestrierung von Multi-Container-Anwendungen. Dabei werden alle Services, 
Netzwerke und Volumes in einer einzigen \texttt{docker-compose.yml}-Datei 
beschrieben und mit einem einzigen Befehl gestartet. Im Kontext dieser 
Diplomarbeit wird Docker eingesetzt, um alle Serverkomponenten -- Backend-Services, 
Datenbanken, API-Gateway und Authentifizierungsdienst -- containerisiert und 
reproduzierbar bereitzustellen.

% -------------------------------------------------------------------
\subsection{Microservices}
% Autor: Niklas | Umfang: 3–4 Seiten
% TODO: Wie Kubernetes funktioniert. Eventuell Unterschied zu Docker. Zusammenspiel von beidem

